\documentclass[journal]{IEEEtran}
\IEEEoverridecommandlockouts

% ================================================================
% Packages
% ================================================================
\usepackage{cite}
\usepackage{amsmath,amssymb,amsfonts}
\usepackage{algorithmic}
\usepackage{graphicx}
\usepackage{textcomp}
\usepackage{xcolor}
\usepackage{booktabs}
\usepackage{multirow}
\usepackage{enumitem}
\usepackage{subcaption}
\usepackage{hyperref}
\usepackage{cleveref}
\usepackage{siunitx}
\usepackage{threeparttable}

% ================================================================
% Math commands
% ================================================================
\newcommand{\MAE}{\mathrm{MAE}}
\newcommand{\Wone}{W_1}
\newcommand{\Xtrain}{X_{\mathrm{train}}}
\newcommand{\Xtest}{X_{\mathrm{test}}}
\newcommand{\Xcore}{X_{\mathrm{core}}}
\newcommand{\Dtrain}{D_{\mathrm{train}}}
\newcommand{\Dtest}{D_{\mathrm{test}}}
\newcommand{\Dcore}{D_{\mathrm{core}}}
\newcommand{\ftrain}{f_{\mathrm{train}}}
\newcommand{\fcore}{f_{\mathrm{core}}}
\newcommand{\Ltrain}{L_{\mathrm{train}}}
\newcommand{\Lcore}{L_{\mathrm{core}}}
\newcommand{\edisc}{\varepsilon_{\mathrm{disc}}}
\newcommand{\Gtest}{G_{\mathrm{test}}}
\newcommand{\Jkmed}{J_{\mathrm{kmed}}}

% Draft mode helpers
\newcommand{\todo}[1]{{\color{red}\textbf{[TODO: #1]}}}
\newcommand{\rewrite}[1]{{\color{blue}#1}}

\begin{document}

% ================================================================
\title{Quantization Cost as a Distribution-Aware Proxy \\
for Coreset Selection in Traffic Forecasting}
% ================================================================

\author{
  \todo{Author names}
  \thanks{\todo{Affiliations and funding}}
}

\maketitle

% ================================================================
\begin{abstract}
% ================================================================
Deep learning models for large-scale traffic forecasting require
substantial computation to train on massive spatio-temporal datasets.
\emph{Coreset selection}---identifying a small, representative subset
of training samples---can reduce this computational burden, but a
critical question remains: how can one predict, before training,
which coreset will produce the best-performing model?

We address this question with three contributions.
First, we conduct a systematic comparison of coreset selection methods
(k-medoids, k-center, graph-cut) across multiple distance metrics and
two forecasting architectures (STGCN, AGCRN) on a real-world traffic
dataset comprising 893~sensors.
Second, we propose \emph{quantization cost}---the average distance
from each training sample to its nearest coreset point---as a
pre-training proxy metric for coreset quality, achieving
within-setting Pearson~$r > 0.94$ with downstream test MAE.
Third, we provide theoretical justification through a Wasserstein-1
bound on the test-set performance gap between full-data and coreset
models.
The bound reveals that the only term controllable by the selection
algorithm is the weighted transport cost
$\Wone(\Xtrain, C; \mathbf{w})$, which k-medoids directly minimises.
Experiments confirm that k-medoids consistently achieves the lowest
test MAE and is the only method for which the coreset distribution
can be closer to the test distribution than the full training set.
\end{abstract}

\begin{IEEEkeywords}
Traffic forecasting, coreset selection, data pruning,
Wasserstein distance, k-medoids, proxy metric
\end{IEEEkeywords}


% ================================================================
\section{Introduction}
\label{sec:introduction}
% ================================================================

\IEEEPARstart{T}{he} advancement of Intelligent Transportation
Systems (ITS) has increased the importance and applicability of
traffic prediction.
Traffic forecasting plays a pivotal role in enhancing road traffic
efficiency, supporting applications such as travel time estimation,
congestion management, route guidance, and signal
control~\cite{vlahogianni_short-term_2014}.
The availability of extensive traffic data from network-wide loop
detectors---for example, the Caltrans Performance Measurement System
(PeMS) provides flow, speed, and density at 5-minute intervals from
thousands of sensors across
California~\cite{chen_freeway_2001}---has accelerated the development
and validation of data-driven prediction models.

Recent innovations in deep learning have demonstrated that Graph
Neural Networks (GNNs) can capture complex spatio-temporal
dependencies in traffic data, achieving superior performance compared
to traditional methods such as ARIMA and
LSTM~\cite{wu_graph_2019, yu_spatio-temporal_2018, cui_traffic_2019}.
However, these models require substantial training data and extensive
computation.
Recent benchmarks show that training state-of-the-art GNN models on
large-scale sensor networks can require over 100~GPU-hours on a single
GPU~\cite{liu_largest_2023}.
As sensor networks grow to thousands of detectors and data accumulates
over months or years, the computational cost of training becomes a
significant barrier to rapid model development and deployment.

\textit{Coreset selection}, the task of identifying a small,
representative subset of training samples from the original dataset,
offers a principled approach to this challenge.
In computer vision, coreset selection has demonstrated significant
advances in image classification: models trained on coresets using
only 50\% of the original data can achieve predictive performance
within a few percent of the full-data
baseline~\cite{xia_moderate_2022, maharana_d2_2023, tan_data_2025}.
In our prior work~\cite{our_itsc}, we presented the first systematic
benchmark of model-free coreset selection methods for GNN-based traffic
forecasting, showing that 30--50\% of the training data is sufficient
to maintain prediction accuracy near the full-data level.
Furthermore, we observed that certain coresets (often at 50--90\%
ratios) slightly \emph{outperformed} models trained on the entire
dataset, suggesting that coreset selection can serve as an implicit
form of data cleaning.

However, two important gaps remain.
First, there is no established \emph{pre-training quality metric}:
given a candidate coreset, practitioners cannot predict how well the
resulting model will perform without actually training it.
Such a metric would enable rapid evaluation of selection algorithms,
distance functions, and subset ratios, significantly reducing the
trial-and-error cost of coreset design.
Second, there is no \emph{theoretical understanding} of why certain
selection methods (notably k-medoids) consistently outperform others
in the traffic forecasting setting.
Without such understanding, practitioners lack principled guidance
for choosing a selection algorithm.

In this paper, we address both gaps.
Building on the empirical foundation of~\cite{our_itsc}, we extend
our investigation from a single model (STGCN) to multiple
architectures, from the PEMS datasets to a different regional
network, and---most importantly---from purely empirical comparison
to a theoretically grounded framework.

The main contributions of this paper can be summarised as follows:
\begin{itemize}
  \item We conduct a systematic empirical comparison of three coreset
    selection methods (k-medoids, k-center, graph-cut) with four
    distance metrics across two forecasting architectures (STGCN,
    AGCRN) on a large-scale traffic dataset (893~sensors, 3~months).
    k-medoids consistently achieves the lowest test MAE across all
    settings.

  \item We propose \emph{quantization cost}---the average distance
    from each training sample to its nearest coreset point---as a
    pre-training proxy metric for coreset quality.
    Within each experimental setting (fixed model and distance type),
    quantization cost achieves Pearson~$r > 0.94$ with downstream
    test MAE, substantially outperforming alternative metrics such as
    Sinkhorn divergence ($r = -0.86$) or effective sample
    size ($r = -0.73$).

  \item We derive a Wasserstein-1 bound on the test-set performance
    gap between the full-data and coreset models.
    The bound identifies the weighted transport cost
    $\Wone(\Xtrain, C; \mathbf{w})$ as the only term directly
    controllable by the selection algorithm, and we prove that
    k-medoids minimises an upper bound on this term.
    We further explain why this bound is particularly tight for
    traffic forecasting, where the MAE loss is Lipschitz and
    targets are near-deterministic.
\end{itemize}

The remainder of this paper is organised as follows.
Section~\ref{sec:related} reviews related work.
Section~\ref{sec:formulation} formalises the problem.
Section~\ref{sec:theory} presents the theoretical analysis.
Section~\ref{sec:framework} describes the proposed framework.
Section~\ref{sec:experiments} reports experimental results.
Section~\ref{sec:conclusion} concludes.


% ================================================================
\section{Related Work}
\label{sec:related}
% ================================================================

This section reviews previous work in three areas relevant to our
study: deep learning for traffic forecasting, coreset selection
methods, and distributional quality metrics.

\subsection{Deep Learning for Traffic Forecasting}
\label{sec:related-traffic}

Traffic forecasting is a multivariate spatio-temporal time-series
problem: given historical observations from a sensor network, the
goal is to predict future traffic states (e.g., flow, speed,
occupancy) across multiple locations.
Since traffic states are influenced by both adjacent locations and
past temporal patterns, it is essential to capture both spatial
correlations among sensors and temporal dependencies across time
steps.

GNNs have emerged as the dominant paradigm for this task, as they
can model road networks explicitly as graphs whose nodes represent
sensors and edges represent spatial
connectivity~\cite{yin_deep_2021}.
Spatio-temporal GNNs have made significant progress in recent
years.
STGCN~\cite{yu_spatio-temporal_2018} combines graph convolutions
with temporal convolutions for efficient feature extraction.
DCRNN~\cite{li_diffusion_2018} employs diffusion convolutions with
a recurrent architecture to model directed traffic flow.
Graph WaveNet~\cite{wu_graph_2019} introduces an adaptive adjacency
matrix alongside dilated causal convolutions, removing the
dependence on a predefined graph structure.
AGCRN~\cite{bai_adaptive_2020} further extends adaptive graph
learning with node-specific parameters and a GRU-based temporal
module.
More recently, Transformer-based architectures have been applied
to traffic forecasting, capturing long-range temporal dependencies
through attention mechanisms~\cite{xu_spatial-temporal_2020,
shao_decoupled_2022, li_dynamic_2024, wang_adaptive_2024}.

Despite these architectural advances, the \emph{data efficiency}
aspect of traffic forecasting has received comparatively little
attention.
Most studies assume that more training data is always beneficial,
yet the computational cost of training scales linearly with dataset
size.
To address these bottlenecks, efficient learning paradigms have been
proposed, including model compression and architectural
optimisation~\cite{zheng_lightcast_2024, zhang_efficient_2025}.
However, these approaches focus on \emph{model} efficiency rather
than \emph{data} efficiency.
Coreset selection offers a complementary, model-agnostic approach
to reducing the computational burden by operating directly on the
training data.


\subsection{Coreset Selection and Data Pruning}
\label{sec:related-coreset}

Coreset selection aims to identify a small but representative subset
of a training dataset while minimising the performance degradation
when a model is trained on the subset rather than the full data.
This technique directly reduces the computational burden associated
with training deep models on massive datasets.
Previous research has studied coreset selection extensively in the
image classification domain, and methods can be broadly categorised
into two groups.

\subsubsection{Model-free methods}
These methods select coresets purely based on data characteristics
without requiring model training or gradient information.
\emph{Geometric methods} exploit distance metrics in feature space.
K-center greedy~\cite{sener_active_2018} iteratively selects the
sample farthest from the current coreset, minimising the covering
radius.
Herding~\cite{welling_herding_2009} constructs a coreset by
iteratively adding samples to match the empirical mean of the
full dataset.
K-medoids~\cite{kaufmann_clustering_1987} partitions data into
clusters and selects actual data points as centres (medoids),
minimising the average distance from each point to its nearest
medoid.
\emph{Submodular methods} leverage submodular objective functions
to measure the information content of a subset.
Facility location~\cite{iyer_submodular_2021,
kothawade_prism_2022} selects samples that maximise total coverage,
while graph-cut maximises cross-partition similarity to encourage
diversity.

\subsubsection{Model-dependent methods}
These methods utilise information from a trained model (e.g.,
gradients, losses, or learned representations) to select coresets.
\emph{Gradient-based methods} select subsets whose aggregated
gradients approximate the full-data gradient:
CRAIG~\cite{mirzasoleiman_coresets_2020} formulates this as a
facility location problem over gradient distances, while
GradMatch~\cite{killamsetty_grad-match_2021} uses orthogonal
matching pursuit.
\emph{Score-based methods} rank samples by importance scores such
as error norms (EL2N), gradient norms
(GraNd)~\cite{paul_deep_2021}, or forgetting
events~\cite{toneva_empirical_2018}.
\emph{Hybrid methods} combine model-derived information with
geometric principles:
Moderate~\cite{xia_moderate_2022} selects samples near the median
distance to class centres,
D2~Pruning~\cite{maharana_d2_2023} balances difficulty and
diversity via message passing, and
Infomax~\cite{tan_data_2025} maximises information while
penalising redundancy.

The vast majority of coreset selection research targets image
classification.
Applying coreset selection to multivariate time-series forecasting,
particularly traffic prediction, has received limited attention.
In our prior work~\cite{our_itsc}, we presented the first benchmark
of eight model-free methods for traffic forecasting, demonstrating
that k-medoids achieves the strongest overall performance by
maintaining broad feature-space coverage.
The present study extends this line of work by providing a
\emph{theoretical} explanation for why k-medoids excels and
proposing a principled proxy metric for coreset quality.

Since our primary motivation is to reduce the computational overhead
of training large-scale forecasting models, we focus on
\emph{model-free} methods that do not require additional model
training during the selection phase.


\subsection{Distributional Quality Metrics}
\label{sec:related-metrics}

Optimal transport (OT) distances provide a natural framework for
measuring distributional similarity between datasets.
The Wasserstein-1 distance ($\Wone$), also known as the earth
mover's distance, computes the minimum cost of transporting mass
from one distribution to another and has strong theoretical
connections to the Kantorovich--Rubinstein
duality~\cite{villani_optimal_2009}.
The Sinkhorn divergence~\cite{cuturi_sinkhorn_2013} offers a
computationally efficient, entropy-regularised approximation.
Maximum Mean Discrepancy (MMD) provides a kernel-based alternative
that is widely used in domain adaptation.

These metrics have been employed in various machine learning
contexts, including dataset distillation, domain adaptation, and
fairness evaluation.
However, their use as \emph{predictors of coreset quality}---that
is, as pre-training indicators of how well a model trained on a
coreset will perform---has not been systematically investigated.
We fill this gap by evaluating multiple distributional metrics as
proxy predictors and demonstrating that quantization cost
(directly connected to $\Wone$) substantially outperforms
alternatives.


% ================================================================
\section{Problem Formulation}
\label{sec:formulation}
% ================================================================

\subsection{Traffic Flow Prediction}

The objective of traffic flow prediction is to forecast future
traffic states based on historical observations over a sensor
network.
Let $\mathbf{X} \in \mathbb{R}^{N \times T \times F}$ denote the
multivariate time-series data collected from $N$ sensors over $T$
time steps with $F$ features (e.g., flow, speed, occupancy).
Given a window of $T_{\mathrm{in}}$ historical observations, the
goal is to predict traffic states for the next
$T_{\mathrm{out}}$ steps:
\begin{equation}\label{eq:forecast}
  \hat{\mathbf{Y}}_{t:t+T_{\mathrm{out}}-1}
  = f_\theta\!\left(
    \mathbf{X}_{t-T_{\mathrm{in}}:t-1}
  \right),
\end{equation}
where $\hat{\mathbf{Y}} \in \mathbb{R}^{N \times T_{\mathrm{out}}
\times 1}$ is the predicted traffic flow sequence.
The parameters $\theta$ are learnt by minimising the mean absolute
error (MAE) between predictions and ground-truth targets
$\mathbf{Y}$:
\begin{equation}\label{eq:mae}
  \MAE_D(f) := \frac{1}{|D|}
  \sum_{(X, Y) \in D} \| f(X) - Y \|_1.
\end{equation}

The dataset $\Dtrain = \{(X_i, Y_i)\}_{i=1}^{n}$ is formed by
sliding a window over the time series, and is partitioned
chronologically into training, validation, and test sets.


\subsection{Coreset Selection}

Given the entire training dataset $\Dtrain$, \emph{coreset
selection} aims to identify a representative subset
$\Dcore \subset \Dtrain$ of size $|\Dcore| = k = \lfloor rn \rfloor$
for a specified selection ratio $r \in (0, 1)$.
Let $\ftrain$ and $\fcore$ denote models trained on $\Dtrain$ and
$\Dcore$, respectively.
The quality of a coreset is measured by the test-set performance gap:
\begin{equation}\label{eq:gtest}
  \Gtest := \bigl|
    \MAE_{\Dtest}(\ftrain) - \MAE_{\Dtest}(\fcore)
  \bigr|.
\end{equation}
The ideal coreset minimises $\Gtest$ while maximising the data
reduction (i.e., minimising $r$).


\subsection{Selection Methods}
\label{sec:methods}

We evaluate three representative model-free coreset selection
methods.
For all methods that require distance calculations, we operate on
PCA-reduced feature representations (Section~\ref{sec:distance}).

\textbf{k-Medoids}~\cite{kaufmann_clustering_1987}.
Partitions the training data into $k$ clusters and selects actual
data points as cluster centres (medoids) by minimising the sum of
distances from each point to its nearest medoid:
\begin{equation}
  \Jkmed(C) = \frac{1}{n}\sum_{i=1}^{n}\min_{c \in C} d(x_i, c).
\end{equation}
We employ the FasterPAM algorithm for scalable computation on
large datasets.

\textbf{k-Center Greedy}~\cite{sener_active_2018}.
Iteratively selects the data point that is farthest from the
current coreset, minimising the maximum distance from any
training point to its nearest coreset member:
$J_{\mathrm{kcen}}(C) = \max_{i}\min_{c \in C} d(x_i, c)$.

\textbf{Graph-Cut}~\cite{iyer_submodular_2021}.
Maximises cross-partition similarity between selected and
unselected points.
Using a Gaussian kernel
$s_{ij} = \exp(-\|x_i - x_j\|^2 / 2\sigma^2)$, the graph-cut
function is
$f_{\mathrm{GC}}(A) = 2\sum_{i \in \Omega} \sum_{a \in A} s_{ia}
- \sum_{a, a' \in A} s_{aa'}$,
encouraging both coverage and diversity.


\subsection{Distance Metrics and Feature Space}
\label{sec:distance}

Computing pairwise distances directly on the raw input features
$\mathbf{X} \in \mathbb{R}^{N \times T_{\mathrm{in}} \times F}$
is computationally expensive and may suffer from the curse of
dimensionality.
To address this, we project each sample into a lower-dimensional
space using Principal Component Analysis (PCA) with
$d_{\mathrm{PCA}} = 10$ components, which captures 82.2\% of the
total variance in our dataset.

Within this PCA embedding, we evaluate four distance types:
\begin{itemize}
  \item \textbf{Euclidean}: distance on the flattened
    $T \times N \times F$ representation.
  \item \textbf{Temporal}: distance on the mean-over-nodes
    representation, capturing temporal patterns.
  \item \textbf{Spatial}: distance on the mean-over-time
    representation, capturing spatial patterns.
  \item \textbf{Combined}: normalised sum of temporal and spatial
    distances, capturing both dimensions simultaneously.
\end{itemize}


% ================================================================
\section{Theoretical Analysis}
\label{sec:theory}
% ================================================================

In this section, we derive a bound on the test-set performance gap
$\Gtest$ and establish that k-medoids is the natural selection
algorithm for minimising this bound.
We then discuss conditions under which a coreset model can
\emph{outperform} the full-data model, and explain why this
theoretical framework is particularly well-suited to traffic
forecasting.

\subsection{Notation and Assumptions}

Let $d$ be a metric on the input space $\mathcal{X}$
(Euclidean distance in PCA space).
For a fixed model~$f$, define the pointwise loss
$\phi_f(x) := |f(x) - y(x)|$,
where $y(x)$ is the ground-truth target associated with input~$x$.
We assume $\phi_f$ is $L_f$-Lipschitz with respect to~$d$, i.e.,
$|\phi_f(x) - \phi_f(x')| \le L_f \, d(x, x')$ for all
$x, x' \in \mathcal{X}$.

\subsection{Performance Gap Bound}

\begin{theorem}\label{thm:main}
Let $C = \{c_1, \ldots, c_k\}$ be the selected coreset points
with cluster-size weights $w_\ell = n_\ell / n$, where $n_\ell$
is the number of training points assigned to medoid $c_\ell$.
Define the weighted coreset MAE as
$\MAE^w_C(f) := \sum_{\ell=1}^{k} w_\ell \, |f(c_\ell) - y(c_\ell)|$
and the model discrepancy as
$\edisc := |\MAE^w_C(\ftrain) - \MAE^w_C(\fcore)|$.
Then:
\begin{equation}\label{eq:main-bound}
\boxed{
  \Gtest \le (\Ltrain + \Lcore)
  \bigl[\Wone(\Xtest, \Xtrain) + \Wone(\Xtrain, C; \mathbf{w})\bigr]
  + \edisc,
}
\end{equation}
where $\Wone(\Xtrain, C; \mathbf{w})$ is the Wasserstein-1 distance
with target marginal~$\mathbf{w}$:
\begin{equation}\label{eq:w1w}
  \Wone(\Xtrain, C; \mathbf{w})
  := \min_{\substack{\Gamma \ge 0 \\
    \sum_\ell \Gamma_{i\ell} = 1/n \\
    \sum_i \Gamma_{i\ell} = w_\ell}}
  \sum_{i,\ell} \Gamma_{i\ell} \, d(x_i, c_\ell).
\end{equation}
\end{theorem}

\begin{proof}
We decompose $\Gtest$ via the triangle inequality:
\begin{align}\label{eq:triangle}
  \Gtest &\le
    \underbrace{|\MAE_{\Dtest}(\ftrain) -
      \MAE_{\Dtrain}(\ftrain)|}_{\text{(A) test--train shift,
      train model}}
    \nonumber \\
  &\quad + \underbrace{|\MAE_{\Dtrain}(\ftrain) -
      \MAE_{\Dtrain}(\fcore)|}_{\text{(B) model gap on train set}}
    \nonumber \\
  &\quad + \underbrace{|\MAE_{\Dtrain}(\fcore) -
      \MAE_{\Dtest}(\fcore)|}_{\text{(C) test--train shift,
      core model}}.
\end{align}

\emph{Step~1 (Bounding A and C).}
Both terms compare a fixed model's MAE on the test set versus the
training set.
By the Kantorovich--Rubinstein (KR) inequality for Lipschitz
functions on empirical measures:
\begin{equation}
  \text{(A)} \le \Ltrain \Wone(\Xtest, \Xtrain), \quad
  \text{(C)} \le \Lcore \Wone(\Xtest, \Xtrain).
\end{equation}

\emph{Step~2 (Bounding B).}
We introduce $\MAE^w_C$ as a waypoint, decomposing term~(B) into
three parts via a second application of the triangle inequality:
\begin{align}
  \text{(B)} &\le
    |\MAE_{\Dtrain}(\ftrain) - \MAE^w_C(\ftrain)|
    + \edisc \nonumber \\
  &\quad + |\MAE^w_C(\fcore) - \MAE_{\Dtrain}(\fcore)|.
\end{align}
The first and third terms compare a fixed model's performance on the
training set versus its weighted evaluation on coreset points.
Applying the weighted KR inequality yields:
\begin{equation}
  \text{(B)} \le (\Ltrain + \Lcore) \Wone(\Xtrain, C; \mathbf{w})
  + \edisc.
\end{equation}

Combining Steps~1 and~2 yields~\eqref{eq:main-bound}.
\end{proof}

\textbf{Interpretation.}
The bound separates $\Gtest$ into three distinct contributions:
\begin{itemize}
  \item $\Wone(\Xtest, \Xtrain)$: the distributional shift between
    test and training sets, determined entirely by the data split.
    This term is \emph{uncontrollable} by the coreset design.

  \item $\Wone(\Xtrain, C; \mathbf{w})$: the distributional mismatch
    between the training set and coreset.
    This is the \emph{only term directly controllable} by the
    selection algorithm.

  \item $\edisc$: the model discrepancy, measuring how differently
    the full-data and coreset models perform on the coreset points
    under cluster weights.
    This term is indirectly influenced by coreset quality.

  \item $\Ltrain, \Lcore$: the Lipschitz constants of the pointwise
    loss for each model, measuring sensitivity to input perturbation.
\end{itemize}


\subsection{Connection to k-Medoids}
\label{sec:theory-kmed}

Having identified $\Wone(\Xtrain, C; \mathbf{w})$ as the
controllable term, we now establish its connection to the k-medoids
objective.

\begin{proposition}\label{prop:kmed}
For any coreset $C = \{c_1, \ldots, c_k\} \subset \Xtrain$ with
nearest-medoid assignment weights~$\mathbf{w}$:
\begin{equation}\label{eq:kmed-w1}
  \Wone(\Xtrain, C; \mathbf{w}) \le \Jkmed(C).
\end{equation}
\end{proposition}

\begin{proof}
The nearest-medoid assignment
$\Gamma_{i\ell} = \frac{1}{n} \mathbf{1}[\pi(x_i) = c_\ell]$,
where $\pi(x_i) = \arg\min_{c \in C} d(x_i, c)$,
defines a feasible transport plan for~\eqref{eq:w1w}:
its row marginals sum to $1/n$ (each point is assigned to exactly
one medoid) and its column marginals sum to $w_\ell = n_\ell / n$
(each medoid receives mass proportional to its cluster size).
The cost of this plan is exactly $\Jkmed(C)$.
Since $\Wone$ is the minimum over all feasible plans,
$\Wone \le \Jkmed$.
\end{proof}

\textbf{Design principle.}
Minimising $\Jkmed$ (the k-medoids objective) directly shrinks the
only controllable term in the bound~\eqref{eq:main-bound}.
Among the three methods we compare---k-medoids, k-center, and
graph-cut---only k-medoids is designed to minimise this quantity.
k-center minimises the \emph{maximum} rather than \emph{average}
distance, and graph-cut maximises cross-partition similarity, neither
of which provides guarantees on $\Wone(\Xtrain, C; \mathbf{w})$.

\begin{remark}[Why not minimise $\Wone$ directly?]
\label{rem:w1-direct}
Computing the exact $\Wone(\Xtrain, C; \mathbf{w})$ requires solving
a linear program of size $O(n \times k)$, which is cubic in the
worst case.
In contrast, k-medoids operates in $O(nk)$ per iteration with
the FasterPAM algorithm and provides $\Jkmed \ge \Wone$ as an upper
bound.
Empirically, $\Jkmed$ is a tight surrogate for~$\Wone$, and the
bound~\eqref{eq:main-bound} is dominated by~$\edisc$ rather than
the $\Wone$ term (Section~\ref{sec:exp-bound}), making exact
$\Wone$ minimisation an unnecessary computational expense.
\end{remark}


\subsection{When Can the Coreset Model Win?}
\label{sec:theory-win}

An intriguing empirical observation from our prior
work~\cite{our_itsc} is that certain coresets can yield models that
\emph{outperform} the full-data baseline.
We now provide a distributional explanation for this phenomenon.

Define the signed performance gap:
\begin{equation}
  \Delta_{\mathrm{test}} :=
    \MAE_{\Dtest}(\fcore) - \MAE_{\Dtest}(\ftrain).
\end{equation}
We seek conditions under which $\Delta_{\mathrm{test}} < 0$, i.e.,
the coreset model beats the full-data model on the test set.

Motivated by the KR inequality, the test error of a model $f$
trained on dataset $D$ can be decomposed as:
\begin{equation}
  \MAE_{\Dtest}(f)
  \approx \MAE_D(f) + L_f \, \Wone(\Xtest, X_D).
\end{equation}
Applying this to both models and subtracting, under the assumptions
that
(i)~both models fit their respective training data equally well
($\MAE_{\Dtrain}(\ftrain) \approx \MAE_{\Dcore}(\fcore)$) and
(ii)~both models have similar sensitivity
($\Ltrain \approx \Lcore \approx L$, as they share the same
architecture and similar training), the training-set terms cancel:
\begin{equation}\label{eq:delta-approx}
\boxed{
  \Delta_{\mathrm{test}}
  \approx L \bigl[
    \Wone(\Xtest, \Xcore) - \Wone(\Xtest, \Xtrain)
  \bigr].
}
\end{equation}

In words: the performance gap is governed by which
distribution---coreset or full training set---is closer to the test
set.
When $\Wone(\Xtest, \Xcore) < \Wone(\Xtest, \Xtrain)$, the gap
shrinks and may even reverse ($\Delta_{\mathrm{test}} < 0$, coreset
model wins).

\textbf{Role of k-medoids.}
By the triangle inequality for~$\Wone$:
\begin{equation}
  \bigl|
    \Wone(\Xtest, \Xcore) - \Wone(\Xtest, \Xtrain)
  \bigr|
  \le \Wone(\Xtrain, \Xcore).
\end{equation}
When k-medoids makes $\Wone(\Xtrain, \Xcore)$ small
(Section~\ref{sec:theory-kmed}), the coreset and training
distributions become nearly interchangeable from the test set's
perspective:
$\Wone(\Xtest, \Xcore) \approx \Wone(\Xtest, \Xtrain)$.
Moreover, because the coreset is a selective subset, it may exclude
training-time outlier patterns (e.g., incidents or holiday effects)
that are absent from the test period, placing $\Xcore$ slightly
closer to $\Xtest$ than $\Xtrain$ is.
This mechanism explains the empirical observation
from~\cite{our_itsc} that k-medoids coresets at 50--90\% ratios
sometimes outperform full-data training.


\subsection{Why Traffic Forecasting Is a Favourable Setting}
\label{sec:theory-traffic}

Our bound is particularly well-suited to traffic forecasting for
several structural reasons that distinguish it from the
classification setting where most coreset literature is focused:

\begin{enumerate}
  \item \textbf{Lipschitz loss.}
    The MAE loss $\phi_f(x) = |f(x) - y(x)|$ is Lipschitz whenever
    $f$ and $y$ are Lipschitz, which holds for smooth traffic
    patterns.
    In contrast, the 0-1 loss used in classification is
    discontinuous and not Lipschitz, rendering KR-based bounds
    inapplicable.

  \item \textbf{Near-deterministic targets.}
    Traffic flow $Y | X$ is nearly deterministic: the same input
    pattern (time of day, day of week, upstream conditions) produces
    a similar flow value.
    This means the conditional noise
    $\Delta_{\mathrm{cond}} \approx 0$, ensuring that the pointwise
    loss $\phi_f$ varies smoothly as a function of~$x$.

  \item \textbf{Density-proportional optimality.}
    In classification, coreset quality hinges on coverage near
    decision boundaries.
    In regression, prediction error is distributed across the
    entire input space proportionally to the data density, which
    is precisely what k-medoids optimises through its
    density-proportional clustering.

  \item \textbf{Natural temporal distribution shift.}
    Traffic data exhibits inherent temporal variation between
    training and test periods (e.g., seasonal changes, evolving
    demand patterns), making the $\Wone(\Xtest, \Xtrain)$ term
    non-trivial and the distributional analysis practically
    meaningful.
\end{enumerate}


% ================================================================
\section{Proposed Framework}
\label{sec:framework}
% ================================================================

\subsection{Coreset Selection Pipeline}

Our framework comprises three stages, illustrated in
Fig.~\ref{fig:pipeline}:

\textbf{Stage~1: Feature extraction and embedding.}
Each training sample $(X_i, Y_i)$ is flattened and projected to a
$d_{\mathrm{PCA}}$-dimensional space via PCA.
We use $d_{\mathrm{PCA}} = 10$ components, capturing 82.2\% of the
total variance.
This dimensionality reduction serves two purposes: it makes
distance computation tractable for large datasets and mitigates the
curse of dimensionality in high-dimensional feature spaces.

\textbf{Stage~2: Offline coreset selection.}
A selection algorithm (k-medoids, k-center, or graph-cut) operates
on the PCA embeddings to choose $k = \lfloor rn \rfloor$ sample
indices.
Selection is performed \emph{once} offline, and the resulting
indices are stored for reuse across multiple model architectures
and hyperparameter configurations.
This amortises the selection cost over all subsequent training runs.

\textbf{Stage~3: Quality prediction and model training.}
Before committing to an expensive training run, the practitioner
computes the quantization cost (Section~\ref{sec:quant-metric})
as a pre-training quality check.
If the quantization cost is acceptably low, the model is trained
on $\Dcore$ using the standard pipeline.

\todo{Fig.~\ref{fig:pipeline}: Pipeline overview diagram}

\subsection{Quantization Cost as Proxy Metric}
\label{sec:quant-metric}

We propose the \emph{quantization cost} as a pre-training proxy
metric for coreset quality:
\begin{equation}\label{eq:quant}
  Q(C) := \frac{1}{n} \sum_{i=1}^{n}
    \min_{c \in C} d_{\mathrm{PCA}}(x_i, c).
\end{equation}
This is precisely the k-medoids objective $\Jkmed(C)$, and by
Proposition~\ref{prop:kmed}, it upper-bounds the controllable term
$\Wone(\Xtrain, C; \mathbf{w})$ in the bound~\eqref{eq:main-bound}.

Quantization cost possesses several desirable properties:
\begin{itemize}
  \item \textbf{Computationally efficient}: $O(nk)$ after PCA
    projection, requiring no model training or gradient computation.
  \item \textbf{Theoretically grounded}: directly connected to the
    Wasserstein-1 distance through the k-medoids link
    (Proposition~\ref{prop:kmed}).
  \item \textbf{Model-agnostic}: depends only on the data geometry,
    not on the forecasting architecture.
  \item \textbf{Universally applicable}: can be computed for any
    selection method, not only k-medoids.
\end{itemize}


% ================================================================
\section{Experiments}
\label{sec:experiments}
% ================================================================

\subsection{Experimental Setup}
\label{sec:exp-setup}

\textbf{Dataset.}
We use traffic flow data from XTraffic~\cite{}, collected from
893~loop detectors in San Bernardino County, California.
We extract three months of data at 5-minute intervals
(24{,}192~time steps in total), partitioned chronologically into
training, validation, and test sets at a 7:1:2 ratio.

\textbf{Models.}
We evaluate two representative spatio-temporal forecasting
architectures:
\textbf{STGCN}~\cite{yu_spatio-temporal_2018}, which combines
Chebyshev graph convolutions with temporal convolutions, and
\textbf{AGCRN}~\cite{bai_adaptive_2020}, which employs adaptive
graph learning with node-specific GRU modules.
Both models use an input horizon of $T_{\mathrm{in}} = 12$ steps
(60~minutes) and an output horizon of $T_{\mathrm{out}} = 12$ steps.

\textbf{Coreset configurations.}
We evaluate three selection methods (k-medoids, k-center, graph-cut)
with four distance types (euclidean, temporal, spatial, combined)
in PCA-10 space, at selection ratios $r \in \{0.3, 0.7\}$ with
seeds 42 and 123.
This yields $3 \times 4 \times 2 \times 2 = 48$ unique coresets,
each trained with two model architectures, producing 96~training
runs in total.

\textbf{Full-data baselines.}
For each model architecture, we train a full-data baseline
($r = 1.0$) to serve as the reference point for computing $\Gtest$.

\textbf{Metrics.}
Performance is measured using MAE, RMSE, and MAPE.
All reported results use the best validation checkpoint.

% --------------------------------------------------
\subsection{Quantization Cost Predicts MAE}
\label{sec:exp-quant}
% --------------------------------------------------

\todo{Table: Method $\times$ distance, quant cost vs MAE at
ratio = 0.3 and 0.7.
Figure: Scatter plot of quant cost vs MAE, coloured by method.
Key numbers: $r > 0.94$ within-setting, $r = 0.983$ global (STGCN).
graph\_cut quant cost 6--7$\times$ higher $\to$ MAE~+2.5.
Comparison with sinkhorn ($-0.86$), voronoi\_cv ($0.76$),
ess\_ratio ($-0.73$).}

% --------------------------------------------------
\subsection{Coreset Advantage ($\Delta \Wone > 0$)}
\label{sec:exp-advantage}
% --------------------------------------------------

\todo{Table: $SW_1$(test, core), $SW_1$(core, train), $\Delta W_1$
by method.
Key: k-medoids is the only method achieving $\Delta W_1 > 0$
(8 out of 96 cases, all k-medoids).
k-medoids $SW_1$(core, train) = 0.17--0.52
(6--17\% of baseline 3.04).
Mechanism: $\Jkmed$ minimisation $\to$ triangle inequality $\to$
coreset $\approx$ train from test's perspective.}

% --------------------------------------------------
\subsection{Bound Verification}
\label{sec:exp-bound}
% --------------------------------------------------

\todo{Table: (A)(B)(C) decomposition with measured values.
$L$ values very small ($\sim$0.001).
$\edisc$ dominates ($\sim$95\% of bound value).
Bound holds: $L_{95}$ $\sim$77\%, $L_{\max}$ $\sim$82\%.
graph\_cut at ratio = 0.3: 0\% (bound too loose).
Discussion of tightness and dominant terms.}

% --------------------------------------------------
\subsection{Ablation Studies}
\label{sec:exp-ablation}
% --------------------------------------------------

\todo{Distance type effect (combined vs.\ single modality).
Ratio sensitivity (0.3 vs.\ 0.7).
Seed stability across configurations.}


% ================================================================
\section{Discussion}
\label{sec:discussion}
% ================================================================

\todo{
\textbf{Dominant term.} $\edisc$ accounts for $\sim$95\% of the
bound, suggesting that tighter bounds could be achieved by directly
modelling model discrepancy. Future work could explore
regularisation strategies that explicitly minimise $\edisc$.

\textbf{Why graph-cut fails.} Graph-cut's objective (cross-partition
similarity) does not control $\Wone(\Xtrain, C; \mathbf{w})$,
leading to $\sim$6--7$\times$ larger quantization cost than
k-medoids at ratio = 0.3 and substantially higher test MAE.

\textbf{Practical implications.} Quantization cost provides a
computationally cheap ($O(nk)$) ``sanity check'' before committing
to expensive training. Practitioners can compute it for multiple
candidate coresets and select the configuration with the lowest
value.

\textbf{Limitations.} Our evaluation uses a single metropolitan
region (San Bernardino) and two model architectures. Future work
should validate findings on additional datasets (e.g., PEMS04,
PEMS07) and architectures (e.g., Transformer-based models).
The PCA-10 embedding, while effective, may not capture all relevant
structure; alternative representations (e.g., learned embeddings)
warrant investigation.
}


% ================================================================
\section{Conclusion}
\label{sec:conclusion}
% ================================================================

We presented a systematic study of coreset selection for traffic
forecasting, extending beyond empirical benchmarking to provide
theoretical foundations.
Our three contributions are:
(1)~a comprehensive comparison across methods, distances, and
architectures, confirming that k-medoids consistently achieves the
lowest test MAE;
(2)~quantization cost as a pre-training proxy metric with
Pearson~$r > 0.94$ within-setting correlation to downstream
performance; and
(3)~a Wasserstein-1 bound that identifies the train--coreset
transport cost as the only controllable term and proves that
k-medoids directly minimises its upper bound.

These results provide practitioners with both a practical tool
(quantization cost for rapid coreset evaluation) and a theoretical
foundation (the $\Wone$-based bound) for principled coreset design
in traffic forecasting.
Future work will extend our evaluation to additional datasets and
model architectures, investigate gradient-based selection methods
as complements to geometric approaches, and explore tighter bounds
that directly address the model discrepancy term~$\edisc$.


% ================================================================
% References
% ================================================================
\bibliographystyle{IEEEtran}
% \bibliography{references/papers}

\begin{thebibliography}{99}

\bibitem{vlahogianni_short-term_2014}
E.~I.~Vlahogianni, M.~G.~Karlaftis, and J.~C.~Golias,
``Short-term traffic forecasting: Where we are and where we're going,''
\emph{Transportation Research Part C: Emerging Technologies},
vol.~43, pp.~3--19, 2014.

\bibitem{chen_freeway_2001}
C.~Chen, K.~Petty, A.~Skabardonis, P.~Varaiya, and Z.~Jia,
``Freeway performance measurement system: Mining loop detector data,''
\emph{Transportation Research Record}, vol.~1748, no.~1,
pp.~96--102, 2001.

\bibitem{yin_deep_2021}
X.~Yin \emph{et~al.},
``Deep learning on traffic prediction: Methods, analysis, and future directions,''
\emph{IEEE Trans. Intell. Transp. Syst.}, vol.~23, no.~6,
pp.~4927--4943, 2021.

\bibitem{yu_spatio-temporal_2018}
B.~Yu, H.~Yin, and Z.~Zhu,
``Spatio-temporal graph convolutional networks: A deep learning framework for traffic forecasting,''
in \emph{Proc. IJCAI}, 2018.

\bibitem{li_diffusion_2018}
Y.~Li, R.~Yu, C.~Shahabi, and Y.~Liu,
``Diffusion convolutional recurrent neural network: Data-driven traffic forecasting,''
in \emph{Proc. ICLR}, 2018.

\bibitem{wu_graph_2019}
Z.~Wu, S.~Pan, G.~Long, J.~Jiang, and C.~Zhang,
``Graph WaveNet for deep spatial-temporal graph modeling,''
\emph{arXiv preprint arXiv:1906.00121}, 2019.

\bibitem{bai_adaptive_2020}
L.~Bai, L.~Yao, C.~Li, X.~Wang, and C.~Wang,
``Adaptive graph convolutional recurrent network for traffic forecasting,''
in \emph{Proc. NeurIPS}, vol.~33, pp.~17804--17815, 2020.

\bibitem{cui_traffic_2019}
Z.~Cui, K.~Henrickson, R.~Ke, and Y.~Wang,
``Traffic graph convolutional recurrent neural network,''
\emph{IEEE Trans. Intell. Transp. Syst.}, vol.~21, no.~11,
pp.~4883--4894, 2019.

\bibitem{xu_spatial-temporal_2020}
M.~Xu \emph{et~al.},
``Spatial-temporal transformer networks for traffic flow forecasting,''
\emph{arXiv preprint arXiv:2001.02908}, 2020.

\bibitem{shao_decoupled_2022}
Z.~Shao \emph{et~al.},
``Decoupled dynamic spatial-temporal graph neural network for traffic forecasting,''
\emph{Proc. VLDB Endowment}, vol.~15, no.~11, pp.~2733--2746, 2022.

\bibitem{li_dynamic_2024}
Z.~Li, J.~Zhou, Z.~Lin, and T.~Zhou,
``Dynamic spatial aware graph transformer for spatiotemporal traffic flow forecasting,''
\emph{Knowledge-Based Systems}, vol.~297, p.~111946, 2024.

\bibitem{wang_adaptive_2024}
R.~Wang \emph{et~al.},
``Adaptive spatio-temporal relation based transformer for traffic flow prediction,''
\emph{IEEE Trans. Vehicular Technology}, 2024.

\bibitem{liu_largest_2023}
X.~Liu \emph{et~al.},
``LargeST: A benchmark dataset for large-scale traffic forecasting,''
in \emph{Proc. NeurIPS}, vol.~36, pp.~75354--75371, 2023.

\bibitem{zheng_lightcast_2024}
R.~Zheng, D.~Zhang, C.~Dong, S.~Huang, and J.~Wang,
``LightCast: Efficient traffic flow forecasting via an integrated compression framework,''
\emph{IEEE Trans. Intelligent Vehicles}, 2024.

\bibitem{zhang_efficient_2025}
Q.~Zhang, X.~Gao, H.~Wang, S.~M.~Yiu, and H.~Yin,
``Efficient traffic prediction through spatio-temporal distillation,''
in \emph{Proc. AAAI}, 2025.

\bibitem{sener_active_2018}
O.~Sener and S.~Savarese,
``Active learning for convolutional neural networks: A core-set approach,''
in \emph{Proc. ICLR}, 2018.

\bibitem{welling_herding_2009}
M.~Welling,
``Herding dynamical weights to learn,''
in \emph{Proc. ICML}, 2009.

\bibitem{kaufmann_clustering_1987}
L.~Kaufmann and P.~Rousseeuw,
``Clustering by means of medoids,''
in \emph{Data Analysis Based on the L1-Norm and Related Methods},
pp.~405--416, 1987.

\bibitem{iyer_submodular_2021}
R.~Iyer, N.~Khargoankar, J.~Bilmes, and H.~Asanani,
``Submodular combinatorial information measures with applications in machine learning,''
in \emph{Proc. ICML}, 2021.

\bibitem{kothawade_prism_2022}
S.~Kothawade, V.~Kaushal, G.~Ramakrishnan, J.~Bilmes, and R.~Iyer,
``PRISM: A rich class of parameterized submodular information measures for guided data subset selection,''
in \emph{Proc. AAAI}, 2022.

\bibitem{mirzasoleiman_coresets_2020}
B.~Mirzasoleiman, J.~Bilmes, and J.~Leskovec,
``Coresets for data-efficient training of machine learning models,''
in \emph{Proc. ICML}, 2020.

\bibitem{killamsetty_grad-match_2021}
K.~Killamsetty, S.~Durga, G.~Ramakrishnan, A.~De, and R.~Iyer,
``Grad-Match: Gradient matching based data subset selection for efficient deep model training,''
in \emph{Proc. ICML}, 2021.

\bibitem{paul_deep_2021}
M.~Paul, S.~Ganguli, and G.~K.~Dziugaite,
``Deep learning on a data diet: Finding important examples early in training,''
in \emph{Proc. NeurIPS}, vol.~34, pp.~20596--20607, 2021.

\bibitem{toneva_empirical_2018}
M.~Toneva, A.~Sordoni, R.~T.~des~Combes, A.~Trischler, Y.~Bengio,
and G.~J.~Gordon,
``An empirical study of example forgetting during deep neural network learning,''
\emph{arXiv preprint arXiv:1812.05159}, 2018.

\bibitem{xia_moderate_2022}
X.~Xia, J.~Liu, J.~Yu, X.~Shen, B.~Han, and T.~Liu,
``Moderate coreset: A universal method of data selection for real-world data-efficient deep learning,''
in \emph{Proc. ICLR}, 2022.

\bibitem{maharana_d2_2023}
A.~Maharana, P.~Yadav, and M.~Bansal,
``D2 Pruning: Message passing for balancing diversity and difficulty in data pruning,''
\emph{arXiv preprint arXiv:2310.07931}, 2023.

\bibitem{tan_data_2025}
H.~Tan, S.~Wu, W.~Huang, S.~Zhao, and X.~Qi,
``Data pruning by information maximization,''
in \emph{Proc. ICLR}, 2025.

\bibitem{sorscher_beyond_2022}
B.~Sorscher, R.~Geirhos, S.~Shekhar, S.~Ganguli, and A.~Morcos,
``Beyond neural scaling laws: Beating power law scaling via data pruning,''
in \emph{Proc. NeurIPS}, 2022.

\bibitem{villani_optimal_2009}
C.~Villani,
\emph{Optimal Transport: Old and New}.
Berlin: Springer, 2009.

\bibitem{cuturi_sinkhorn_2013}
M.~Cuturi,
``Sinkhorn distances: Lightspeed computation of optimal transport,''
in \emph{Proc. NeurIPS}, vol.~26, 2013.

\bibitem{our_itsc}
\todo{Our ITSC paper reference.}

\end{thebibliography}

\end{document}
